\documentclass[12pt]{article}
\usepackage[T1]{fontenc}
\usepackage[utf8]{inputenc}
\usepackage{lmodern}
\usepackage{textcomp}
\usepackage{enumitem}
\usepackage{geometry}
\geometry{margin=1in}
\begin{document}
\begin{center}
Answer Key - Chemistry - Graduate Level Assessment 17 \
\small (Answers are ordered exactly as in the question paper.)
\end{center}

\section*{Multiple Choice}
\begin{enumerate}[label=\arabic*.]
\item \textbf{Answer:} B
\\ \textit{Options:} A) hydrogen bonds; B) free radicals; C) allotropes; D) isotopes; E) nonpolar molecules
\\ \textit{Note:} Free radicals are defined as species that contain unpaired electrons in their Lewis structures, which makes them highly reactive and unstable.
\item \textbf{Answer:} A
\\ \textit{Options:} A) O$_{2}$; B) NO 3; C) NO 2-; D) CO 32-; E) NO 2
\\ \textit{Note:} O$_{2}$ has the shortest bonds due to its double bond character between the carbon and oxygen atoms, resulting in stronger and shorter bond lengths compared to the single bonds found in CO 32-, NO 2-, and NO 3-.
\item \textbf{Answer:} A
\\ \textit{Options:} A) 310.2 MHz; B) 500.5 MHz; C) 275.0 MHz; D) 400.0 MHz; E) 120.2 MHz
\\ \textit{Note:} The NMR frequency of 31 P in a magnetic field is determined using the formula ν = γB 0/(2π), where γ (the gyromagnetic ratio for 31 P) is approximately 17.235 MHz/T, and substituting B$_{0}$ = 20.0 T yields a frequency of approximately 310.2 MHz.
\item \textbf{Answer:} A
\\ \textit{Options:} A) very electronegative; B) typically metallic; C) highly colored; D) mostly gases at room temperature; E) often very reactive with water
\\ \textit{Note:} Monatomic ions of the representative elements tend to be very electronegative because they typically form by gaining or losing electrons to achieve a stable electron configuration, with nonmetals in particular having a strong tendency to attract electrons to complete their valence shell.
\item \textbf{Answer:} A
\\ \textit{Options:} A) H$_{2}$Se; B) H$_{2}$CO$_{3}$; C) H$_{2}$S; D) HBr; E) H$_{2}$O
\\ \textit{Note:} H$_{2}$Se is the strongest acid among the options because it has a larger atomic radius and lower electronegativity compared to other hydrogen chalcogenides, resulting in weaker H-Se bonds that break more easily, thus facilitating proton donation.
\end{enumerate}

\section*{True / False}
\begin{enumerate}[label=\arabic*.]
\setcounter{enumi}{5}
\item \textbf{Answer:} False
\\ \textit{Options:} A) True; B) False
\\ \textit{Note:} The statement is false because the energy of yellow light with a wavelength of 5890 Å is approximately 2.11 eV, not 1.5 eV, which corresponds to about 49.5 Kcal/mole, indicating a discrepancy in both energy values.
\item \textbf{Answer:} False
\\ \textit{Options:} A) True; B) False
\\ \textit{Note:} The swimmer's body temperature does not increase due to the warm breeze after exiting the cold water because the body's heat loss mechanisms, such as evaporative cooling from wet skin, typically dominate in such conditions, leading to a decrease in body temperature rather than an increase.
\item \textbf{Answer:} False
\\ \textit{Options:} A) True; B) False
\\ \textit{Note:} The statement is false because when MgBr 2 and Na$_{2}$CO$_{3}$ are mixed in solution, they undergo a double displacement reaction that results in the formation of insoluble MgCO 3, leading to the precipitation of products and leaving no significant ionic species in solution.
\item \textbf{Answer:} False
\\ \textit{Options:} A) True; B) False
\\ \textit{Note:} The angular-momentum quantum number \( l \) for a \( t \) orbital, which corresponds to the \( d \) subshell, can only be 2, as \( l \) values are defined such that \( l = 0 \) for \( s \), \( l = 1 \) for \( p \), \( l = 2 \) for \( d \), and \( l = 3 \) for \( f \) orbitals.
\item \textbf{Answer:} True
\\ \textit{Options:} A) True; B) False
\\ \textit{Note:} The statement is true because the thermodynamic equilibrium constant K for a reaction at equilibrium is defined as the ratio of the forward rate constant (k\_a) to the reverse rate constant (k\_b), reflecting the relationship between the rates of the forward and reverse reactions.
\end{enumerate}

\section*{Long Form Response}
\begin{enumerate}[label=\arabic*.]
\setcounter{enumi}{10}
\item \textbf{Answer:} The limiting high-temperature molar heat capacity at constant volume (C\_V) of a gas-phase diatomic molecule approaches 4.5 R due to contributions from translational, rotational, and vibrational degrees of freedom. At high temperatures, the translational motion contributes 3 R, while rotational motion adds 2 R (1 for each of the two rotational axes), resulting in a total of 5 R. However, vibrational modes, which begin to contribute significantly at elevated temperatures, introduce an additional factor of 0.5 R per mode due to quantum effects, leading to a total of 4.5 R when considering the vibrational modes are partially excited. Thus, the limiting high-temperature C\_V reflects the balance of these energy contributions, following the equipartition theorem, underlining the significance of molecular structure in thermodynamic behavior.
\\ \textit{Note:} correct because, at high temperatures, the total molar heat capacity for a diatomic gas includes contributions from translational (3 R), rotational (2 R), and vibrational motions, but the vibrational contributions do not fully come into play until temperatures are sufficiently high, leading the limiting value of C\_V to approach 4.5 R instead of 5 R.
\item \textbf{Answer:} The de Broglie wavelength of a proton is given by the equation \( \lambda = \frac{h}{p} \), where \( h \) is Planck's constant and \( p \) is the momentum of the proton. To achieve a wavelength of \( 1.0 \times 10^{-10} \, \text{m} \), the proton must fall through a potential difference that results in the necessary kinetic energy to reach that momentum. The kinetic energy \( K \) gained by the proton when falling through a potential \( V \) is given by \( K = qV \), where \( q \) is the charge of the proton. By equating the kinetic energy to the momentum relation, and solving, we find that the potential required for the proton to achieve the specified wavelength is approximately \( 0.030 \, \text{V} \). This relationship illustrates the quantum mechanical principles governing particle behavior and energy transformation in electric fields.
\\ \textit{Note:} correct because it accurately links the de Broglie wavelength to the proton's momentum, which is derived from its kinetic energy gained through falling through a potential difference, thereby showing how both concepts are interrelated in determining the required conditions for achieving the specified wavelength.
\end{enumerate}

\vfill
\noindent\textit{End of Answer Key}
\end{document}